\section{Performance Analysis of \glitter}
\label{sec:performance}




\begin{figure}[t]
  \centering
        \begin{tikzpicture}
                \begin{groupplot}[group style = {group size = 5 by 1, horizontal sep = 35pt}, width = 8.0cm, height = 7.0cm]
                        \nextgroupplot[ title = {},
                        xmax=29.9,
                        ymin=0.2,
                        xlabel={\textsf{Number of parties ($|C| = n$)}},
                        ylabel={\textsf{Wall clock time (ms)}},
                        ymode=log,
                        grid=major,
                        legend columns=1,
                        legend style={nodes={scale=0.575},draw=black, fill=white, text=black, anchor=north west,at={(1,1)}},
                        ]

                        \newcommand\tline[2]{
                        \addplot [color=#1, mark=*,mark options=solid,mark size=1.25pt,
                            error bars/.cd, y dir=both, y explicit,
                            error bar style={color=black,line width=0.75pt,solid},
                            error mark options={line width=1pt,mark size=1.75pt,rotate=90}]
                            table [x=n, y=tot_ms, y error=totsd]{t#2.dat}
                            node[right,pos=1,font=\scriptsize] {$t=#2$};
                            }

                        \tline{red}{3}
                        \tline{orange}{5}
                        \tline{olive}{6}
                        \tline{blue}{7}
                        \tline{purple}{8}
                        \tline{brown}{9}
                        \tline{black}{11}

                        \addplot [color=black!80!white, mark=+,mark options=solid,mark size=1.25pt,
                            error bars/.cd, y dir=both, y explicit,
                            error bar style={color=black,line width=0.75pt,solid},
                            error mark options={line width=1pt,mark size=1.75pt,rotate=90}]
                            table [x=n, y=tot_ms]{bullet.dat}
                            node[below,pos=0.1,font=\scriptsize] {MuSig-DN};
                \end{groupplot}
        \end{tikzpicture}
\caption{Single-core wall clock time for for \glitter.
%The times shown are the sum of the computation times for
%$\sign_1$, $\sign_2$, and $\combine$. The computation time for MuSig-DN
%is shown for comparison.
}
\label{fig:exp-nt}
\end{figure}

\begin{figure}[t]
  \centering
%\documentclass{standalone}
\usepackage{pgfplots}
\usepackage{relsize}
\usepackage{versions}
\usetikzlibrary{pgfplots.groupplots}
\pgfplotsset{compat=1.3}

\begin{document}
        \begin{tikzpicture}
                \begin{groupplot}[group style = {group size = 5 by 1, horizontal sep = 35pt}, width = 5.0cm, height = 6.0cm]
                        \nextgroupplot[ title = {},
                        xtick={1,2,4,8,16,32},
                        xticklabels={1,2,4,8,16,32},
                        xlabel={\textsf{Number of cores}},
                        ylabel={\textsf{Wall clock time (ms)}},
                        xmode=log,
                        ymode=log,
                        grid=major,
                        legend columns=1,
                        legend style={nodes={scale=0.575},draw=black, fill=white, text=black, anchor=north west,at={(1,1)}},
                        ]

                        \newcommand\speedupline[3]{
                        \addplot [color=#1, mark=*,mark options=solid,mark size=1.25pt,
                            error bars/.cd, y dir=both, y explicit,
                            error bar style={color=black,line width=0.75pt,solid},
                            error mark options={line width=1pt,mark size=1.75pt,rotate=90}]
                            table [x=nthreads, y=tot_ms, y error=totsd]{s#2.dat}
                            %node[left,pos=0,font=\scriptsize] {#3}
                            node[left,pos=1,font=\sffamily\tiny] {#3}
                            ;
                            }

                        \speedupline{black}{25}{40-core 2.3\,GHz~~~}
                        \speedupline{black}{25.squish}{4-core 3.7\,GHz}
                        %\speedupline{black}{23}{$n=23,t=11$}{??$\times$ speedup}
                        %\speedupline{black}{21}{$n=21,t=11$}{??$\times$ speedup}

                \end{groupplot}
        \end{tikzpicture}
\end{document}

        \begin{tikzpicture}
                \begin{groupplot}[group style = {group size = 5 by 1, horizontal sep = 35pt}, width = 6.0cm, height = 7.0cm]
                        \nextgroupplot[ title = {},
                        xtick={1,2,4,8,16,32},
                        xticklabels={1,2,4,8,16,32},
                        xlabel={\textsf{Number of cores}},
                        ylabel={\textsf{Wall clock time (ms)}},
                        xmode=log,
                        ymode=log,
                        grid=major,
                        legend columns=1,
                        legend style={nodes={scale=0.575},draw=black, fill=white, text=black, anchor=north west,at={(1,1)}},
                        ]

                        \newcommand\speedupline[3]{
                        \addplot [color=#1, mark=*,mark options=solid,mark size=1.25pt,
                            error bars/.cd, y dir=both, y explicit,
                            error bar style={color=black,line width=0.75pt,solid},
                            error mark options={line width=1pt,mark size=1.75pt,rotate=90}]
                            table [x=nthreads, y=tot_ms, y error=totsd]{s#2.dat}
                            %node[left,pos=0,font=\scriptsize] {#3}
                            node[left,pos=1,font=\sffamily\tiny] {#3}
                            ;
                            }

                        \speedupline{black}{25}{40-core 2.3\,GHz~~~}
                        \speedupline{black}{25.squish}{4-core 3.7\,GHz}
                        %\speedupline{black}{23}{$n=23,t=11$}{??$\times$ speedup}
                        %\speedupline{black}{21}{$n=21,t=11$}{??$\times$ speedup}

                \end{groupplot}
        \end{tikzpicture}
\caption{Scaling experiment, showing the wall clock time of $(n,|\coalition|,t) = (25, 25, 11)$ shown in Figure~\ref{fig:exp-nt} relative to CPU cores.
  }
\label{fig:exp-scaling}
\end{figure}


In this section, we analyze the performance of \glitter.  In terms of
the number of rounds of communication, \glitter matches the state of the
art, with two, and it sends significantly less bandwidth per signature,
at 65 bytes per participant.  Therefore, we focus on the computational
complexity.

There are two sources of potentially expensive computation: the two
calls to $\DVRFOnem.\generate$ (one in each of $\sign_1$ and $\sign_2$),
and the call to $\DVRFOnem.\verify$ in $\sign_2$.  Which one dominates
depends on the parameters $n$, $t$, and the number of participants in the signing protocol.
Recall that $\coalition$ is a set representing the identifiers of participants in a particular signing session.
Although the minimum number of participants required for signing is
$\lvert \coalition \rvert \geq \minparticipants \geq 2t-1$,
we assume $\lvert \coalition \rvert = n$ for this analysis,
to give an upper performance bound.
As such, performance will be even better when $\minparticipants \leq \lvert \coalition \rvert < n$.

$\DVRFOnem.\generate$ primarily performs $\delta = {n-1 \choose t-1}$
hash computations, field multiplications, and field additions, as seen
in Equation~\ref{eq:vrf-out-eval} (recalling that the $\cramerpoly_{a_i}(\partyid)$ values can be precomputed).
$\DVRFOnem.\verify$ computes $\ell=|\coalition|-t$,
$|\coalition|$-way multiexponentiations.  When $t$ is small, we expect
the $\DVRFOnem.\verify$ cost to dominate, and for larger $t$, the
$\DVRFOnem.\generate$ cost should dominate.

To concretely evaluate the performance of \glitter, we implemented it in
Rust.\footnote{Our code is available at
\url{https://git-crysp.uwaterloo.ca/iang/arctic/}}.
We ran our implementation over all allowable combinations of
parameters $t \ge 2$, $2t-1 \le |\coalition| \le n \le 25$, using a
% this is squish
4-core 3.7\,GHz Intel E-2374G CPU.
We measured the computation time for each of $\sign_1$,
$\sign_2$, and $\combine$, averaged over 10 signatures, for each parameter
combination.  In Figure~\ref{fig:exp-nt}, we show the (single-core)
total runtime of $\sign_1$, $\sign_2$, and $\combine$ for various
combinations of parameters.  We concretely measure $\DVRFOnem.\generate$
to take around $0.24 \delta$ microseconds for each of its two
invocations, and $\DVRFOnem.\verify$ to take around $7 \ell
|\coalition|$ microseconds.  For $t \le 4$, the latter dominates, for
$t=5$, they are roughly comparable, and for $t \ge 6$, the former
quickly dominates.

For comparison, we also show the computational time for MuSig-DN, but
\emph{only} the zero-knowledge proof and verification components of
their algorithm.  We ran their code~\cite{musig-dn-code} on our same
machine to obtain these figures.  We can see that for $n\le 20$,
\glitter is more than an order of magnitude faster than MuSig-DN, and
for $n \le 10$, it is three orders of magnitude faster.

As seen in Figure~\ref{fig:exp-nt}, the computation time for the $(n,|\coalition|,t)
= (25, 25, 11)$ parameter combination, where $\delta = {24 \choose 10} =
1961256$ is around 940\,ms, almost all of which is spent in
$\DVRFOnem.\generate$ computing Equation~\ref{eq:vrf-out-eval}.
However, we observe that Equation~\ref{eq:vrf-out-eval} computes the sum of $\delta$
independent terms, and so is highly amenable to parallelization, which we also implemented and measured.
%As this CPU only had four
%cores, and so the parallelization potential was limited,
We ran this scaling experiment both on the above machine, and also on
% this is sizzle
a 40-core 2.3\,GHz Intel 8380 CPU.  The results are shown in
Figure~\ref{fig:exp-scaling}.  Although the slower clock speed of the
40-core CPU puts it at a disadvantage for smaller numbers of cores,
we can see almost linear scaling for both CPUs.  Using all 4 cores, the
4-core CPU sees a speedup of 3.69$\times$ for $\DVRFOnem.\generate$ and
3.65$\times$ in total time, while using 32 cores, the 40-core CPU sees a
speedup of 27.9$\times$ for $\DVRFOnem.\generate$ and 25.7$\times$ in total time.
Beyond 32 cores, we observed diminishing returns.

%\ian{Perhaps some discussion of applications for \glitter where either
%$n$ and $t$ are small, and so it's already fast, and/or applications
%where $n$ and $t$ are larger, but operations happen at a human scale (a
%human needs to manually approve a message to be signed), so a one-second
%processing time is no problem?  You seemed to have applications of
%\glitter in mind in an earlier conversation?  Can we put them here?}
%\chelsea{i don't want to make up applications, and i don't have anything that
%concrete in my head, beyond wallets.}


